\section{Study Design}\label{sec:study-design}

We design a controlled experiment to assess coding agent capabilities on
real feature implementation tasks drawn from production repositories.
This section describes the task selection, agent configurations, prompt
design, evaluation framework, and experiment execution protocol.

%% ------------------------------------------------------------------
\subsection{Task Selection}\label{sec:tasks}

We select 12 tasks from four open-source repositories that are actively
maintained by Huawei and the broader community.  Each task corresponds
to a merged pull request or commit that implements a complete feature or
bug fix.  Table~\ref{tab:tasks} summarises the task registry.

\begin{table}[t]
  \caption{Task registry.  GT = ground truth (merged code).  Complexity
    is classified by file count, cross-component coordination required,
    and domain-specific knowledge.}
  \label{tab:tasks}
  \centering
  \small
  \begin{tabular}{llllrr}
    \toprule
    \textbf{ID} & \textbf{Repository} & \textbf{Complexity} & \textbf{Lang} & \textbf{GT Files} & \textbf{GT Lines} \\
    \midrule
    C1 & cann-ops-adv  & Low    & C++          &   1 &       3 \\
    C2 & cann-ops      & Low    & C++          &   4 &       9 \\
    M1 & MindSpeed     & Low    & Python       &   3 &      19 \\
    K1 & kubernetes    & Low    & Go           &  13 &   1{,}010 \\
    \midrule
    C3 & torch\_npu    & Medium & Python       &   9 &     637 \\
    C4 & cann-ops      & Medium & C++/AscendC  &  24 &   1{,}273 \\
    M2 & MindSpeed     & Medium & Python       &   6 &     150 \\
    K2 & kubernetes    & Medium & Go           &  35 &   3{,}759 \\
    K3 & kubernetes    & Medium & Go           &  49 &  11{,}000 \\
    \midrule
    C5 & cann-ops      & High   & C++/AscendC  &  27 &   3{,}372 \\
    M3 & MindSpeed     & High   & Python       &  10 &     948 \\
    K4 & kubernetes    & High   & Go           &  98 &   6{,}507 \\
    \bottomrule
  \end{tabular}
\end{table}

\textbf{Selection criteria.}
Every task was drawn from a real, merged contribution to a production
repository.  CANN tasks (C1--C5) and MindSpeed tasks (M1--M3) come from
Huawei's Ascend NPU software stack---operator kernels, training
framework features, and precision fixes that represent typical
industrial workloads.  Kubernetes tasks (K1--K4) come from the
Kubernetes project and cover KEP-driven feature gates, API types,
controllers, and authorization changes, representing large-scale
open-source engineering.

\textbf{Complexity classification.}
We classify tasks into three tiers based on the number of ground-truth
files, cross-component coordination requirements, and domain-specific
knowledge needed:
\begin{itemize}
  \item \emph{Low} (C1, C2, M1, K1): 1--13 files, localised changes
    within a single module or subsystem.
  \item \emph{Medium} (C3, C4, M2, K2, K3): 6--49 files, requiring
    coordination across multiple components (e.g., host and device code,
    API types and controllers).
  \item \emph{High} (C5, M3, K4): 10--98 files, demanding architectural
    understanding, multi-module orchestration, and extensive testing.
\end{itemize}

The tasks span four languages (C++, AscendC, Python, Go) and three
complexity tiers, providing a diverse substrate for evaluating agent
capabilities across different engineering contexts.

%% ------------------------------------------------------------------
\subsection{Agent Configurations}\label{sec:configs}

We evaluate four primary agent configurations, summarised in
Table~\ref{tab:configs}.  Each configuration pairs a frontier foundation model
with an agent framework.

\begin{table}[t]
  \caption{Agent configurations evaluated on the full 12-task benchmark.
    Each runs all $12 \times 2 = 24$ task--prompt combinations.}
  \label{tab:configs}
  \centering
  \small
  \begin{tabular}{lllr}
    \toprule
    \textbf{Config} & \textbf{Model} & \textbf{Framework} & \textbf{$N$} \\
    \midrule
    Claude  & Claude Opus 4.6  & Claude Code & 24 \\
    Codex   & GPT-5.4          & Codex CLI   & 24 \\
    Cursor  & Composer-2       & Cursor Agent & 24 \\
    OpenCode & GLM-5.1         & OpenCode    & 24 \\
    \bottomrule
  \end{tabular}
\end{table}

\textbf{Product-level comparisons.}
Each configuration pairs a different model with a different agent framework,
so comparisons confound model capability with framework engineering.
These comparisons are best understood as \emph{product-level}
assessments---comparing the capabilities of commercially available agent
tools---rather than controlled model ablations.
We discuss this confound further in \S\ref{sec:threats}.

%% ------------------------------------------------------------------
\subsection{Prompt Design}\label{sec:prompts}

Each task is paired with two prompt variants that differ in the level
of specification detail provided to the agent:

\begin{itemize}
  \item \textbf{Short prompt}: contains a \emph{Summary} (one-sentence
    problem statement) and a \emph{Proposal} (one-sentence solution
    direction).  Typically 2--4 sentences.
  \item \textbf{Long prompt}: extends the short prompt with
    \emph{Motivation} (why the change is needed) and \emph{Design
    Details} (numbered list of implementation considerations).
    Typically 15--30 sentences.
\end{itemize}

\textbf{Design principle.}
Both prompt levels describe \emph{what} to implement and \emph{why},
but never reveal specific file paths, function names, or code
snippets from the ground truth.  This ensures that agents must
independently navigate the codebase to locate relevant files and design
their implementation.

\textbf{Example (C1, buffer alias fix).}
The short prompt reads: ``\emph{ScaledMaskedSoftmaxGradV2 has a tensor
lifetime violation in ComputeSoftmaxGrad---a buffer is used after its
queue resource is freed.  Replace the reference with an equivalent
surviving buffer.}''  The long prompt adds three numbered design
details: root cause analysis of the data flow, the fix strategy
(identify a shape/dtype-compatible live buffer), and the impact scope
(NormHeadDim variant only).

This two-level design allows us to measure the \emph{prompt
granularity effect}: the extent to which richer specification
compensates for agent limitations in codebase understanding and
architectural reasoning.

%% ------------------------------------------------------------------
\subsection{Evaluation Framework}\label{sec:eval-framework}

\textbf{Ground truth.}
For each task, the ground truth is the actual merged code from the
repository (commit or pull request).  This provides a concrete, human-reviewed
reference implementation against which we compare agent output.

\textbf{Scoring dimensions.}
Each experiment is scored by an LLM evaluator on three dimensions, each
on a 0--5 scale:
\begin{itemize}
  \item \textbf{Score~A (Functional Correctness)}: Does the generated
    code correctly implement the intended functionality?
  \item \textbf{Score~B (Completeness \& Coverage)}: Does the
    implementation cover all required files and components?
  \item \textbf{Score~C (Behavioral Equivalence)}: Does the
    implementation match the ground truth's design decisions, API
    conventions, and architectural approach?
\end{itemize}

\textbf{Verdict rules.}
We assign a three-level verdict based on the dimension scores:
\begin{itemize}
  \item \textsc{pass}: $A \geq 4$ \textbf{and} $B \geq 4$ \textbf{and} $C \geq 3$.
  \item \textsc{fail}: $A < 2$ \textbf{or} $B < 2$ \textbf{or} $C < 2$ (any dimension critically deficient).
  \item \textsc{partial}: all other cases.
\end{itemize}

\textbf{File coverage metric.}
We additionally compute file-level overlap between the agent's generated
files and the ground-truth file set:
$\mathit{coverage} = |\mathit{generated} \cap \mathit{GT}| \;/\;
|\mathit{GT}|$.  This provides a coarse quantitative complement to the
expert-scored dimensions.

\textbf{Evaluation scope.}
Our evaluation measures \emph{production-level equivalence} to the merged implementation, not purely functional correctness.
An agent solution that correctly addresses the requirement but uses different file organization, naming conventions, or architectural patterns than the ground truth will score lower on Behavioral Equivalence~(C) and file coverage.
This design choice reflects our industrial interest: we assess whether agent output could replace the human-authored PR, not merely whether it compiles and passes tests.
As a consequence, our \textsc{fail} verdicts should be read as ``not production-equivalent'' rather than ``functionally incorrect.''
Our root-cause analysis (\S\ref{sec:root-cause}) confirms this: the majority of non-PASS cases stem from engineering completeness gaps, not semantic misunderstanding.

\textbf{Multi-evaluator cross-validation.}
To mitigate evaluator bias, every experiment is independently scored by
three LLM evaluators from different model families: Claude Opus~4.6,
Codex GPT-5.4, and GLM-5.1.  We report the \emph{majority verdict}
(agreement of $\geq$2/3 evaluators) and scores averaged across all
three.  Unanimous agreement is 62\%, and $\geq$2/3 consensus reaches
99\%; pairwise agreement ranges from 69\% to 79\%, with per-dimension
score spreads averaging $<$1~point.

%% ------------------------------------------------------------------
\subsection{Experiment Execution}\label{sec:execution}

\textbf{Sanitization protocol.}
To prevent agents from extracting ground-truth code from repository
history, we apply a strict sanitization procedure to each experiment
repository:
\begin{enumerate}
  \item The repository is rebased to a single commit immediately before
    the target feature, removing all commits that contain ground-truth
    code.
  \item All git remotes are stripped and the \texttt{gh} CLI is blocked,
    preventing agents from fetching upstream history or pull request
    contents.
  \item Agents retain access to the full codebase (for context) and to
    web search (for documentation), but cannot access the specific PR or
    commit that constitutes the ground truth.
\end{enumerate}

\textbf{Anti-cheating audit.}
During evaluation of initial Codex runs on task~K4, we discovered that
the agent had executed \texttt{git log} and \texttt{git cherry-pick}
commands to extract ground-truth commits from an insufficiently
sanitized repository.  The K4 repository was re-sanitized (all target
commits removed from history) and the affected experiments were re-run.
Post-sanitization results were confirmed genuine through inspection of
agent trajectories and comparison of generated file counts and
implementation details against ground truth.
This incident underscores a critical finding: frontier models actively
search for shortcuts in version control history, making thorough
sanitization essential for any benchmark using real repositories.

\textbf{Experiment scale.}
In total, we executed 96 experiments across 12~tasks $\times$
4~configurations $\times$ 2~prompt levels---a complete factorial design
with no missing cells.
Each task--config--prompt combination was run exactly once and
evaluated by three independent LLM judges.
The full results are reported in \S\ref{sec:results}.
